\chapter{Metodologia}

\section{Classificação da Pesquisa}
A presente pesquisa classifica-se quanto à sua natureza como \textbf{aplicada}, pois objetiva gerar conhecimento para aplicação prática dirigida à solução de problemas específicos na robótica autônoma. Quanto aos objetivos, enquadra-se como \textbf{exploratória e experimental}, visto que envolve a manipulação de variáveis (arquiteturas de hardware e modelos de redes neurais) para analisar o comportamento do objeto de estudo (o braço robótico) em condições controladas. A abordagem é mista (qualitativa e quantitativa), combinando o desenvolvimento de software com a análise estatística de desempenho e latência \cite{gil2002}.

\section{Materiais e Instrumentos}
Para a realização dos experimentos, serão utilizados os seguintes recursos, baseados na infraestrutura desenvolvida no projeto de iniciação científica anterior \cite{gic2025}:
\begin{itemize}
    \item \textbf{Hardware Mecânico:} Manipulador robótico antropomórfico "Atom" (baseado na plataforma \textit{InMoov}), manufaturado em PLA via impressão 3D.
    \item \textbf{Hardware de Processamento:}
    \begin{itemize}
        \item \textit{Microcontroladores:} Arduino Mega 2560 (controle de atuadores) e ESP32-CAM (aquisição de imagem sem fio).
        \item \textit{Single Board Computer (SBC):} Raspberry Pi 4 ou Banana Pi M2 Zero (processamento de IA na borda).
    \end{itemize}
    \item \textbf{Software e Ferramentas:} Linguagens Python e C++, bibliotecas OpenCV 4.x, TensorFlow Lite (para inferência embarcada), ambiente de desenvolvimento VS Code e ferramentas de CAD para eventuais ajustes mecânicos.
\end{itemize}

\section{Procedimentos Metodológicos}
O desenvolvimento do projeto seguirá um ciclo incremental dividido em quatro fases distintas:

\subsection{Fase 1: Arquitetura de Visão e Aquisição de Dados (O Olho)}
Nesta etapa, será configurado o sistema de percepção visual. O módulo ESP32-CAM será programado para atuar como um nó de captura de vídeo IP (\textit{stream}), transmitindo quadros em resolução QVGA/VGA via protocolo UDP/TCP para a unidade de processamento central (SBC).
Serão aplicadas técnicas de pré-processamento digital de imagens utilizando a biblioteca OpenCV, incluindo:
\begin{itemize}
    \item Conversão de espaço de cores (BGR para RGB/Grayscale);
    \item Filtros de redução de ruído (Gaussian Blur) para estabilizar a detecção em ambientes com iluminação variável;
    \item Redimensionamento de quadros (\textit{resizing}) para adequação à entrada da rede neural.
\end{itemize}

\subsection{Fase 2: Inteligência Artificial e Processamento na Borda (O Cérebro)}
Esta fase concentra o núcleo da inovação do projeto, focando na implementação de \textit{Edge AI}:
\begin{enumerate}
    \item \textbf{Coleta e Curadoria de Dataset:} Criação de um banco de imagens proprietário contendo os objetos de manipulação (ex: cubos, ferramentas), com anotação de \textit{bounding boxes} para treinamento supervisionado.
    \item \textbf{Treinamento e Otimização:} Utilização de \textit{Transfer Learning} em arquiteturas leves como \textbf{MobileNetV2-SSD} ou \textbf{YOLOv8 Nano}. O treinamento será realizado em ambiente de alto desempenho (PC com GPU) e, posteriormente, o modelo será convertido para o formato \texttt{.tflite} (TensorFlow Lite), aplicando quantização (INT8) para reduzir o uso de memória na SBC.
    \item \textbf{Inferência Embarcada:} Implementação do algoritmo de detecção rodando localmente na Raspberry Pi/Banana Pi, visando uma taxa de atualização mínima de 10 FPS para garantir fluidez.
\end{enumerate}

\subsection{Fase 3: Integração e Controle Cinemático (O Braço)}
A integração entre a "visão" e a "ação" ocorrerá através de protocolos de comunicação serial (UART/I2C):
\begin{itemize}
    \item O algoritmo de IA extrairá as coordenadas do centroide ($X, Y$) e a profundidade estimada ($Z$) do objeto detectado.
    \item Um módulo de cinemática inversa calculará os ângulos necessários para as juntas ($\theta_1, \theta_2, \theta_3, \dots$) para posicionar o efetuador final nas coordenadas alvo.
    \item O Arduino receberá o vetor de ângulos e acionará os servomotores de forma sincronizada.
\end{itemize}

\subsection{Fase 4: Validação e Análise de Dados}
A validação será realizada através de testes de "Pick and Place" (pegar e soltar), coletando métricas quantitativas:
\begin{itemize}
    \item \textbf{Acurácia de Detecção (\%):} Taxa de identificação correta dos objetos em diferentes condições de luz.
    \item \textbf{Taxa de Sucesso na Manipulação (\%):} Percentual de tentativas em que o robô consegue efetivamente segurar e mover o objeto.
    \item \textbf{Latência do Sistema (ms):} Tempo decorrido entre a captura da imagem e o início do movimento do motor.
\end{itemize}

\section{Aspectos Éticos}
Por tratar-se de pesquisa experimental em robótica e engenharia de software, o projeto não envolve testes com seres humanos ou animais, dispensando submissão ao Comitê de Ética em Pesquisa (CEP/CONEP), conforme Resolução CNS nº 510/2016. A coleta de dados visuais restringir-se-á a objetos inanimados no ambiente laboratorial, garantindo que nenhuma imagem de terceiros seja armazenada ou processada sem consentimento.